\documentclass[11pt,a4paper]{article}
% =====================================================================
%  Gemeinsame Vorlage der Arbeitsblätter "Graphen" – Informatik 11 (NTG)
%  (gleiches Layout wie informatik/10/java/arbeitsblaetter/ab-vorlage.tex)
%  Einbinden in jedem Blatt direkt nach \documentclass: % =====================================================================
%  Gemeinsame Vorlage der Arbeitsblätter "Graphen" – Informatik 11 (NTG)
%  (gleiches Layout wie informatik/10/java/arbeitsblaetter/ab-vorlage.tex)
%  Einbinden in jedem Blatt direkt nach \documentclass: % =====================================================================
%  Gemeinsame Vorlage der Arbeitsblätter "Graphen" – Informatik 11 (NTG)
%  (gleiches Layout wie informatik/10/java/arbeitsblaetter/ab-vorlage.tex)
%  Einbinden in jedem Blatt direkt nach \documentclass: \input{ab-vorlage}
%  Übersetzen mit XeLaTeX, LuaLaTeX, pdfLaTeX oder Tectonic.
%
%  Blatt:   xelatex ab-3-1-grundbegriffe.tex
%  Lösung:  xelatex -jobname=ab-3-1-grundbegriffe-lsg "\def\mitloesung{}\input{ab-3-1-grundbegriffe}"
%           (füllt alle Lücken rot aus; siehe \lsg, \lsglinien, \lsgoder)
% =====================================================================
\usepackage[a4paper,left=18mm,right=18mm,top=26mm,bottom=17mm,headheight=15mm,headsep=4mm,footskip=8mm]{geometry}
\usepackage{iftex}
\ifPDFTeX
  \usepackage[T1]{fontenc}
  \usepackage[utf8]{inputenc}
  \usepackage{helvet}
  \usepackage[scaled=0.86]{beramono}
  \renewcommand{\familydefault}{\sfdefault}
\else
  \usepackage{fontspec}
  \setmainfont{texgyreheros}[Extension=.otf,UprightFont=*-regular,BoldFont=*-bold,ItalicFont=*-italic,BoldItalicFont=*-bolditalic]
  \setmonofont{DejaVuSansMono}[Extension=.ttf,UprightFont=*,BoldFont=*-Bold,ItalicFont=*-Oblique,BoldItalicFont=*-BoldOblique,Scale=0.86]
\fi
\usepackage[ngerman,shorthands=off]{babel}
\usepackage{xcolor,graphicx,tabularx,array,enumitem,fancyhdr,tikz,listings,multicol,calc,etoolbox}
\usepackage[most]{tcolorbox}
\usetikzlibrary{arrows.meta,positioning,calc,decorations.pathreplacing,shapes.geometric}
\usepackage[hidelinks]{hyperref}
\setlength{\parindent}{0pt}
\setlength{\parskip}{3pt}
\setlist{nosep,leftmargin=*}

% ---------- Lösungsschalter ----------
\newif\ifloesung
\ifdefined\mitloesung\loesungtrue\fi
\newcommand{\lsgfarbe}{\color{red!75!black}}

% ---------- Kopf- und Fußzeile (einheitliche Form) ----------
\newcommand{\lehrkraft}{Hau}
\newcommand{\themenbereich}{3. Graphen}
\newcommand{\abnummer}{}
\pagestyle{fancy}\fancyhf{}
\renewcommand{\headrulewidth}{0.4pt}
\fancyhead[L]{\small Klasse 11\\Datum:}
\fancyhead[C]{\small Themenbereich:\\\themenbereich}
\fancyhead[R]{\small \lehrkraft\ifloesung\\{\lsgfarbe\bfseries Lösung}\fi}
\fancyfoot[C]{\scriptsize edu-mrh.de\,·\,Informatik 11\,·\,Blatt \abnummer\ifloesung\ (Lösung)\fi\,·\,Seite \thepage}
\newcommand{\abtitel}[2]{\renewcommand{\abnummer}{#1}\begin{center}\Large\bfseries #1\quad\underline{#2}\end{center}\vspace{-1mm}}

% ---------- Boxen ----------
\newenvironment{aufgabe}[2][]{\begin{tcolorbox}[enhanced,breakable,colback=white,colframe=black,boxrule=0.7pt,arc=3mm,left=2mm,right=2mm,top=1.5mm,bottom=1.5mm,before skip=2.5mm,after skip=2.5mm]\textbf{Aufgabe #2)}\ifstrempty{#1}{}{ \textit{#1}}\par\vspace{0.6mm}}{\end{tcolorbox}}
\newtcolorbox{merke}[1][Merke]{enhanced,breakable,colback=green!5,colframe=green!40!black,boxrule=0.8pt,arc=2mm,left=2mm,right=2mm,top=1.5mm,bottom=1.5mm,before skip=2.5mm,after skip=2.5mm,title={#1},fonttitle=\bfseries,coltitle=white}
\newtcolorbox{hinweis}{enhanced,breakable,colback=yellow!12,colframe=orange!70!black,boxrule=0.6pt,arc=2mm,left=2mm,right=2mm,top=1mm,bottom=1mm,before skip=2mm,after skip=2mm}
\newtcolorbox{ausblick}[1][Ausblick]{enhanced,breakable,colback=teal!5,colframe=teal!60!black,boxrule=0.8pt,arc=2mm,left=2mm,right=2mm,top=1.5mm,bottom=1.5mm,before skip=2.5mm,after skip=2.5mm,title={#1},fonttitle=\bfseries,coltitle=white}
\newcommand{\web}[1]{{\small\color{gray}\textit{Website: #1}}}

% ---------- Lücken, Linien, Karos ----------
\newcommand{\luecke}[1][3.5cm]{\rule[-0.5ex]{#1}{0.4pt}}
\newcommand{\linien}[1]{\par\vspace{1.5mm}\foreach \i in {1,...,#1}{\noindent\rule[0pt]{\linewidth}{0.3pt}\par\vspace{5.5mm}}\vspace{-3mm}}
\newcommand{\kariert}[1]{\par\vspace{1.5mm}\noindent\begin{tikzpicture}\draw[step=5mm,gray!55,thin] (0,0) grid (\linewidth-1pt,#1*5mm);\end{tikzpicture}\par}
\newcommand{\karobox}[2]{\begin{tikzpicture}\draw[step=5mm,gray!55,thin] (0,0) grid (#1,#2);\end{tikzpicture}}
\newcommand{\code}[1]{\texttt{#1}}

% Lücke mit Lösung: \lsg[Mindestbreite]{Lösungstext}
\newlength{\lsgbreite}
\newcommand{\lsg}[2][3.5cm]{\ifloesung\settowidth{\lsgbreite}{\,#2\,}\ifdim\lsgbreite<#1\setlength{\lsgbreite}{#1}\fi\underline{\makebox[\lsgbreite][l]{\lsgfarbe\,#2}}\else\luecke[#1]\fi}
% Schreiblinien mit Lösung: \lsglinien{Anzahl}{Lösungstext} (gleiche Höhe wie \linien)
\newcommand{\lsglinien}[2]{\ifloesung\par\vspace{1.5mm}\noindent\begin{minipage}[t][#1\dimexpr 6.9mm\relax][t]{\linewidth}\lsgfarbe #2\end{minipage}\par\vspace{-1.5mm}\else\linien{#1}\fi}
% Beliebiger Inhalt: \lsgoder{nur in der Lösung}{nur im Blatt}
\newcommand{\lsgoder}[2]{\ifloesung#1\else#2\fi}

% ---------- Graphen zeichnen ----------
\tikzset{
  knoten/.style={draw,rounded corners=3pt,fill=orange!12,minimum height=6mm,inner sep=3pt,font=\small},
  kknoten/.style={draw,circle,fill=orange!12,minimum size=7mm,inner sep=1pt,font=\small},
  kante/.style={thick},
  gkante/.style={thick,-{Stealth[length=2.6mm]}},
  dkante/.style={thick,{Stealth[length=2.6mm]}-{Stealth[length=2.6mm]}},
  gewicht/.style={font=\scriptsize,fill=white,inner sep=1pt,text=teal!60!black},
  lsgkante/.style={thick,red!75!black},
  lsggkante/.style={thick,red!75!black,-{Stealth[length=2.6mm]}},
  lsgknoten/.style={knoten,draw=red!75!black,text=red!75!black,fill=red!4},
}
% Adjazenzmatrix zum Ausfüllen: \adjmx[Lösungseinträge]{Anzahl}{Knotenliste}{Zellbreite}
% Lösungseinträge als Zeile/Spalte/Text, z. B. \adjmx[1/2/X,2/1/X]{3}{A,B,C}{7mm}
\newcommand{\adjmx}[4][]{\begin{tikzpicture}[x=#4,y=#4]
  \foreach \n [count=\i] in {#3} {\node[font=\small\bfseries] at (\i,0) {\n}; \node[font=\small\bfseries,anchor=east] at (0.45,-\i) {\n};}
  \foreach \k in {0,...,#2} {\draw[gray] (\k+0.5,-0.5) -- (\k+0.5,-#2-0.5); \draw[gray] (0.5,-\k-0.5) -- (#2+0.5,-\k-0.5);}
  \ifloesung\ifstrempty{#1}{}{\foreach \z/\s/\t in {#1} {\node[text=red!75!black,font=\small] at (\s,-\z) {\t};}}\fi
\end{tikzpicture}}

% ---------- Java-Listings ----------
\lstset{language=Java,basicstyle=\ttfamily\small,keywordstyle=\bfseries,commentstyle=\color{gray!80!black},stringstyle=\color{black},showstringspaces=false,columns=flexible,keepspaces=true,tabsize=3,frame=single,framerule=0.4pt,rulecolor=\color{black},xleftmargin=2mm,framexleftmargin=1mm,aboveskip=2mm,belowskip=2mm,breaklines=true,escapechar=§,
 literate={ä}{{\"a}}1 {ö}{{\"o}}1 {ü}{{\"u}}1 {Ä}{{\"A}}1 {Ö}{{\"O}}1 {Ü}{{\"U}}1 {ß}{{\ss}}1 {–}{{--}}1 {„}{{\glqq}}1 {“}{{\grqq}}1}
\lstnewenvironment{java}[1][]{\lstset{#1}}{}

%  Übersetzen mit XeLaTeX, LuaLaTeX, pdfLaTeX oder Tectonic.
%
%  Blatt:   xelatex ab-3-1-grundbegriffe.tex
%  Lösung:  xelatex -jobname=ab-3-1-grundbegriffe-lsg "\def\mitloesung{}\documentclass[11pt,a4paper]{article}
\input{ab-vorlage}
% Dörfer-Graph (Merkkasten Begriffe): Koordinaten wie im Begriffe-Explorer der Website
\newcommand{\doerferknoten}{%
  \node[knoten,font=\tiny,minimum height=4mm,inner sep=1.5pt] (A) at (4.4,3.3) {Altdorf};
  \node[knoten,font=\tiny,minimum height=4mm,inner sep=1.5pt] (W) at (1.2,2.9) {Weiler};
  \node[knoten,font=\tiny,minimum height=4mm,inner sep=1.5pt] (F) at (5.1,2.0) {Fischbach};
  \node[knoten,font=\tiny,minimum height=4mm,inner sep=1.5pt] (Z) at (2.9,2.0) {Ziegelstein};
  \node[knoten,font=\tiny,minimum height=4mm,inner sep=1.5pt] (N) at (0.9,0.95) {Neustadt};
  \node[knoten,font=\tiny,minimum height=4mm,inner sep=1.5pt] (B) at (4.4,0.8) {Burg};
  \node[knoten,font=\tiny,minimum height=4mm,inner sep=1.5pt] (R) at (2.6,0.2) {Rain};}
\begin{document}
\abtitel{3.1}{Graphen -- Grundbegriffe}

Hänsel und Gretel haben die Hexe überlistet. Doch die Brotkrumen, mit denen sie ihren Weg markiert hatten, haben die Vögel gefressen. Findest du trotzdem nach Hause?

\begin{aufgabe}[Hänsel und Gretel]{1}
Starte das Spiel auf der Website (Station 3.1).
\textbf{a)} Zeichne beim Spielen eine Karte des Waldes: jeder Ort ein Kästchen, jeder Weg eine Linie. Führt ein Weg nur in eine Richtung, zeichne einen Pfeil. Das Hexenhaus ist schon eingezeichnet.
\end{aufgabe}
\noindent\begin{tikzpicture}
\draw[step=5mm,gray!55,thin] (0,0) grid (17.4,7);
\node[knoten,fill=green!15] (hex) at (4,2.3) {Hexenhaus (Start)};
\ifloesung
  \node[lsgknoten] (warn) at (5.5,6.5) {Warnung};
  \node[lsgknoten] (wolf) at (10.5,6.5) {Werwolf};
  \node[lsgknoten] (stadt) at (1.6,5.2) {Stadt};
  \node[lsgknoten] (abz) at (5.5,5.2) {Abzweigung};
  \node[lsgknoten] (spinne) at (13,5.2) {Spinne};
  \node[lsgknoten] (baeck) at (1.6,3.6) {Bäckerei};
  \node[lsgknoten] (pilze) at (4,3.9) {Pilze};
  \node[lsgknoten] (w3) at (13,3.4) {Wald3};
  \node[lsgknoten] (w1) at (8,1.0) {Wald1};
  \node[lsgknoten] (stein) at (12.2,2.0) {Steinkreis};
  \node[lsgknoten] (w2) at (15,0.8) {Wald2};
  \node[lsgknoten] (heim) at (16.2,2.6) {Zuhause};
  \draw[lsgkante] (stadt) -- (baeck);
  \draw[lsgkante] (stadt) -- (pilze);
  \draw[lsgkante] (pilze) -- (abz);
  \draw[lsgkante] (pilze) -- (hex);
  \draw[lsggkante] (pilze) to[bend left=50] node[right,font=\scriptsize,text=red!75!black] {Pilze essen} (hex);
  \draw[lsgkante] (abz) -- (warn);
  \draw[lsggkante] (warn) -- (wolf);
  \draw[lsgkante] (abz) -- (spinne);
  \draw[lsgkante] (spinne) -- (w3);
  \draw[lsgkante] (hex) -- (w1);
  \draw[lsgkante] (w1) -- (w3);
  \draw[lsgkante] (w1) -- (w2);
  \draw[lsgkante] (w1) -- (stein);
  \draw[lsgkante] (w2) -- (w3);
  \draw[lsgkante] (w2) -- (stein);
  \draw[lsgkante] (w3) -- (stein);
  \draw[lsggkante] (w2) -- (heim);
  \node[text=red!75!black,font=\scriptsize,align=left,anchor=north west] at (0.2,1.3) {13 Knoten, 17 Kanten\\(14 ungerichtet, 3 gerichtet)};
\fi
\end{tikzpicture}

\textbf{b)} Am Ende zeigt dir das Spiel deinen Weg. Notiere ihn und die Anzahl der Schritte:
\lsglinien{1}{individuell; kürzester Weg: Hexenhaus -- Wald1 -- Wald2 -- Zuhause (3 Schritte)}

\begin{minipage}[t]{0.6\linewidth}
\begin{merke}[Graph und Kurzschreibweise $G = (V, E)$]\raggedright
Ein Graph besteht aus \lsg[2.3cm]{Knoten} (engl. \textit{vertices}) und \lsg[2.3cm]{Kanten} (engl. \textit{edges}). Eine Kante verbindet zwei Knoten. Im Spiel sind die Orte die \lsg[2cm]{Knoten}, die Wege die \lsg[2cm]{Kanten}.\\[1mm]
Kurzschreibweise $G = (V, E)$: $V$ ist die Menge der Knoten, $E$ die Menge der Kanten.
\begin{itemize}
\item $(A, B)$: \lsg[2.5cm]{gerichtete} Kante, nur von $A$ nach $B$
\item $\{A, B\}$: \lsg[2.5cm]{ungerichtete} Kante, steht für $(A, B)$ und \lsg[1.4cm]{$(B, A)$}
\end{itemize}
\end{merke}
\end{minipage}\hfill
\begin{minipage}[t]{0.37\linewidth}
\vspace{1mm}
\centering\small \textbf{Beispiel: Wer folgt wem?}\\[1mm]
\begin{tikzpicture}[x=0.95cm,y=0.9cm]
\node[knoten] (leon) at (0.7,2.2) {Leon};
\node[knoten] (letifa) at (3.3,2.4) {Letifa};
\node[knoten] (luise) at (4.3,1.1) {Luise};
\node[knoten] (hamza) at (2.4,0.3) {Hamza};
\node[knoten] (sabine) at (0.5,0.6) {Sabine};
\draw[dkante] (leon) -- (letifa);
\draw[gkante] (letifa) -- (luise);
\draw[gkante] (leon) -- (luise);
\draw[dkante] (leon) -- (hamza);
\draw[gkante] (hamza) -- (luise);
\end{tikzpicture}
\end{minipage}\\[1mm]
{\small $V = \{\text{Hamza}, \text{Leon}, \text{Letifa}, \text{Luise}, \text{Sabine}\}$\\
$E = \{\{\text{Leon}, \text{Letifa}\}, \{\text{Leon}, \text{Hamza}\}, (\text{Letifa}, \text{Luise}), (\text{Leon}, \text{Luise}), (\text{Hamza}, \text{Luise})\}$}

\begin{aufgabe}[Hänsel und Gretel]{1}
\textbf{c)} Schreibe den Ausschnitt deiner Karte mit den Knoten Hexenhaus, Pilze, Stadt, Bäckerei und Wald1 in Kurzschreibweise auf (5 Kanten).\\[2mm]
$V = \{$\lsg[15.4cm]{Hexenhaus, Pilze, Stadt, Bäckerei, Wald1}$\}$\\[3mm]
$E = \{$\lsg[15.4cm]{\{Hexenhaus, Pilze\}, (Pilze, Hexenhaus), \{Pilze, Stadt\}, \{Stadt, Bäckerei\}, \{Hexenhaus, Wald1\}}$\}$
\end{aufgabe}

\newpage
\begin{minipage}[t]{0.71\linewidth}
\begin{merke}[Pfade, Zyklen, Richtung, Gewicht]\raggedright\small
\textbf{Pfad (Weg):} Folge von Knoten, in der aufeinanderfolgende Knoten durch eine \lsg[1.8cm]{Kante} verbunden sind; keine Kante wird doppelt benutzt.\\[0.6mm]
\textbf{Einfacher Pfad:} Pfad, in dem sich kein \lsg[1.8cm]{Knoten} wiederholt.\\
Beispiel in A: \lsg[5.6cm]{Altdorf -- Weiler -- Ziegelstein -- Neustadt}\\[0.6mm]
\textbf{Länge eines Pfades:} ungewichtet die \lsg[3.2cm]{Anzahl der Kanten}, gewichtet die \lsg[3.2cm]{Summe der Gewichte}.\\[0.6mm]
\textbf{Zyklus:} einfacher Pfad, dessen \lsg[5cm]{Start- und Endknoten gleich sind}.\\
Beispiel in A: \lsg[5.6cm]{Weiler -- Ziegelstein -- Fischbach -- Weiler}\\
Graph mit Zyklus: \emph{zyklisch}, sonst \lsg[3.4cm]{azyklisch (zyklenfrei)}.\\[0.6mm]
\textbf{Gerichtet:} Kanten nur \lsg[3cm]{in Pfeilrichtung} (Pfeil).\ \textbf{Ungerichtet:} \lsg[2.8cm]{in beide Richtungen} (Strich).\\[0.6mm]
\textbf{Stark zusammenhängend:} Von \lsg[2.4cm]{jedem Knoten} gibt es einen Pfad zu \lsg[3.2cm]{jedem anderen Knoten}.\\
\textbf{Schwach zusammenhängend:} \lsg[5.4cm]{nur ohne Beachtung der Richtungen}\\ \lsg[\linewidth]{hängt alles zusammen. B ist schwach, denn von Rain führt keine Kante weg.}\\[0.6mm]
\textbf{Gewichtet:} Jede Kante trägt eine \lsg[2.6cm]{Zahl (Gewicht)}, z.\,B. \lsg[3.6cm]{Entfernung, Fahrzeit}.
\end{merke}
\end{minipage}\hfill
\begin{minipage}[t]{0.27\linewidth}
\vspace{1mm}
\textbf{\small A} {\scriptsize ungerichtet}\\
\begin{tikzpicture}[x=0.8cm,y=0.85cm]
\doerferknoten
\draw[kante] (A)--(W); \draw[kante] (F)--(A); \draw[kante] (W)--(F); \draw[kante] (W)--(Z);
\draw[kante] (Z)--(F); \draw[kante] (Z)--(N); \draw[kante] (B)--(Z); \draw[kante] (F)--(B);
\draw[kante] (N)--(R); \draw[kante] (B)--(R);
\end{tikzpicture}\\[2mm]
\textbf{\small B} {\scriptsize gerichtet}\\
\begin{tikzpicture}[x=0.8cm,y=0.85cm]
\doerferknoten
\draw[gkante] (A)--(W); \draw[gkante] (F)--(A); \draw[gkante] (W)--(F); \draw[gkante] (W)--(Z);
\draw[gkante] (Z)--(F); \draw[gkante] (Z)--(N); \draw[gkante] (B)--(Z); \draw[gkante] (F)--(B);
\draw[gkante] (N)--(R); \draw[gkante] (B)--(R);
\end{tikzpicture}
\end{minipage}

\begin{aufgabe}{2}
{\raggedright Zeichne $G = (V, E)$ mit $V = \{A, B, C, D, E, F\}$ und\\ $E = \{(A, B), (A, C), (A, E), (B, C), (C, C), (C, D), (E, C), (F, D)\}$. Beantworte die Fragen.\par}\vspace{1mm}
\begin{minipage}[t]{0.42\linewidth}
\vspace{0pt}
\begin{tikzpicture}
\draw[step=5mm,gray!55,thin] (0,0) grid (7,3);
\ifloesung
  \foreach \n/\x/\y/\t in {a/0.6/1.5/A, b/2/2.55/B, e/2/0.45/E, c/3.4/1.5/C, d/4.9/1.5/D, f/6.4/1.5/F} {
    \node[kknoten,minimum size=6mm,draw=red!75!black,text=red!75!black] (\n) at (\x,\y) {\t};}
  \draw[lsggkante] (a)--(b); \draw[lsggkante] (a)--(c); \draw[lsggkante] (a)--(e);
  \draw[lsggkante] (b)--(c); \draw[lsggkante] (c)--(d); \draw[lsggkante] (e)--(c);
  \draw[lsggkante] (f)--(d);
  \draw[lsggkante] (c) to[out=60,in=120,looseness=6] (c);
\fi
\end{tikzpicture}
\end{minipage}\hfill
\begin{minipage}[t]{0.55\linewidth}
\vspace{0pt}\small\raggedright
Gerichtet? Woran erkennbar? \lsg[3.9cm]{ja, Paare ( ) statt \{ \}}\\[1.5mm]
Ein Zyklus: \lsg[5.7cm]{(C, C), eine Schlinge (Länge 1)}\\[1.5mm]
Stark/schwach zusammenhängend? \lsg[3.6cm]{schwach: von D}\\[1.5mm]
\lsg[0.98\linewidth]{kommt man nirgends hin.}\\[1.5mm]
Länge des längsten einfachen Pfades: \lsg[3.2cm]{3 (A--B--C--D)}
\end{minipage}
\end{aufgabe}

\begin{aufgabe}[Zurück in den Wald]{3}
Deine Karte aus Aufgabe 1: \textbf{a)} gerichtet?\ \textbf{b)} ein Zyklus?\ \textbf{c)} stark oder schwach zusammenhängend?\ \textbf{d)} „Der Weg zurück … scheint verschwunden“: Ist die Kante Hexenhaus -- Wald1 gerichtet? Prüfe im Spiel.
\lsglinien{2}{\small a) gemischt; gerichtet nur (Pilze, Hexenhaus), (Warnung, Werwolf), (Wald2, Zuhause)\quad b) Wald1 -- Wald2 -- Wald3 -- Wald1\\ c) schwach: Von Zuhause/Werwolf führt keine Kante weg.\quad d) nein: Von Wald2 aus führt ein Weg zurück.}
\end{aufgabe}

\begin{aufgabe}[Ordne die Graphen ein.]{4}
\vspace{-1mm}
\small\renewcommand{\arraystretch}{1.15}
\begin{tabularx}{\linewidth}{|>{\centering\arraybackslash}m{5.6cm}|>{\centering\arraybackslash}m{1.7cm}|>{\centering\arraybackslash}m{1.1cm}|>{\centering\arraybackslash}m{1.1cm}|>{\centering\arraybackslash}m{1.1cm}|>{\centering\arraybackslash}X|}\hline
\textbf{Graph} & \scriptsize\textbf{zusammen-\newline hängend?}\newline stark/schwach/nein & \scriptsize\textbf{ge-\newline richtet?} & \scriptsize\textbf{ge-\newline wichtet?} & \scriptsize\textbf{zyklisch?} & \scriptsize\textbf{Länge eines längsten einfachen Pfades}\\\hline
\rule{0pt}{8mm}\textbf{I}\ \begin{tikzpicture}[baseline=(k5.base),x=0.85cm,y=0.55cm]
\foreach \n/\x/\y in {1/0/1,2/1/1,3/2/1,4/3/1,5/1/0,6/3/0} {\node[kknoten,minimum size=5mm,inner sep=0pt,font=\tiny] (k\n) at (\x,\y) {K\n};}
\draw[kante] (k1)--(k2)--(k3)--(k4); \draw[kante] (k2)--(k5);
\end{tikzpicture}
& \lsgoder{\lsgfarbe nein}{} & \lsgoder{\lsgfarbe nein}{} & \lsgoder{\lsgfarbe nein}{} & \lsgoder{\lsgfarbe nein}{} & \lsgoder{\lsgfarbe 3: K1--K2--K3--K4}{}\\\hline
\rule{0pt}{5mm}\textbf{II}\ $V = \{A, B, C\}$\newline $E =\{(A, B), (B, A), (C, B), (C, C)\}$
& \lsgoder{\lsgfarbe schwach}{} & \lsgoder{\lsgfarbe ja}{} & \lsgoder{\lsgfarbe nein}{} & \lsgoder{\lsgfarbe ja}{} & \lsgoder{\lsgfarbe 2: C$\to$B$\to$A}{}\\\hline
\rule{0pt}{8mm}\textbf{III}\ \begin{tikzpicture}[baseline=(k3.base),x=0.9cm,y=0.6cm]
\foreach \n/\x/\y in {1/0/1,2/2/1,3/1/0,4/3/0} {\node[kknoten,minimum size=5mm,inner sep=0pt,font=\tiny] (k\n) at (\x,\y) {K\n};}
\draw[kante] (k1) -- node[gewicht] {4} (k2); \draw[kante] (k2) -- node[gewicht] {1} (k3);
\draw[kante] (k3) -- node[gewicht] {5} (k1); \draw[kante] (k3) -- node[gewicht] {2} (k4);
\end{tikzpicture}
& \lsgoder{\lsgfarbe stark}{} & \lsgoder{\lsgfarbe nein}{} & \lsgoder{\lsgfarbe ja}{} & \lsgoder{\lsgfarbe ja}{} & \lsgoder{\lsgfarbe 11: K2--K1--K3--K4}{}\\\hline
\end{tabularx}
\end{aufgabe}
\end{document}
"
%           (füllt alle Lücken rot aus; siehe \lsg, \lsglinien, \lsgoder)
% =====================================================================
\usepackage[a4paper,left=18mm,right=18mm,top=26mm,bottom=17mm,headheight=15mm,headsep=4mm,footskip=8mm]{geometry}
\usepackage{iftex}
\ifPDFTeX
  \usepackage[T1]{fontenc}
  \usepackage[utf8]{inputenc}
  \usepackage{helvet}
  \usepackage[scaled=0.86]{beramono}
  \renewcommand{\familydefault}{\sfdefault}
\else
  \usepackage{fontspec}
  \setmainfont{texgyreheros}[Extension=.otf,UprightFont=*-regular,BoldFont=*-bold,ItalicFont=*-italic,BoldItalicFont=*-bolditalic]
  \setmonofont{DejaVuSansMono}[Extension=.ttf,UprightFont=*,BoldFont=*-Bold,ItalicFont=*-Oblique,BoldItalicFont=*-BoldOblique,Scale=0.86]
\fi
\usepackage[ngerman,shorthands=off]{babel}
\usepackage{xcolor,graphicx,tabularx,array,enumitem,fancyhdr,tikz,listings,multicol,calc,etoolbox}
\usepackage[most]{tcolorbox}
\usetikzlibrary{arrows.meta,positioning,calc,decorations.pathreplacing,shapes.geometric}
\usepackage[hidelinks]{hyperref}
\setlength{\parindent}{0pt}
\setlength{\parskip}{3pt}
\setlist{nosep,leftmargin=*}

% ---------- Lösungsschalter ----------
\newif\ifloesung
\ifdefined\mitloesung\loesungtrue\fi
\newcommand{\lsgfarbe}{\color{red!75!black}}

% ---------- Kopf- und Fußzeile (einheitliche Form) ----------
\newcommand{\lehrkraft}{Hau}
\newcommand{\themenbereich}{3. Graphen}
\newcommand{\abnummer}{}
\pagestyle{fancy}\fancyhf{}
\renewcommand{\headrulewidth}{0.4pt}
\fancyhead[L]{\small Klasse 11\\Datum:}
\fancyhead[C]{\small Themenbereich:\\\themenbereich}
\fancyhead[R]{\small \lehrkraft\ifloesung\\{\lsgfarbe\bfseries Lösung}\fi}
\fancyfoot[C]{\scriptsize edu-mrh.de\,·\,Informatik 11\,·\,Blatt \abnummer\ifloesung\ (Lösung)\fi\,·\,Seite \thepage}
\newcommand{\abtitel}[2]{\renewcommand{\abnummer}{#1}\begin{center}\Large\bfseries #1\quad\underline{#2}\end{center}\vspace{-1mm}}

% ---------- Boxen ----------
\newenvironment{aufgabe}[2][]{\begin{tcolorbox}[enhanced,breakable,colback=white,colframe=black,boxrule=0.7pt,arc=3mm,left=2mm,right=2mm,top=1.5mm,bottom=1.5mm,before skip=2.5mm,after skip=2.5mm]\textbf{Aufgabe #2)}\ifstrempty{#1}{}{ \textit{#1}}\par\vspace{0.6mm}}{\end{tcolorbox}}
\newtcolorbox{merke}[1][Merke]{enhanced,breakable,colback=green!5,colframe=green!40!black,boxrule=0.8pt,arc=2mm,left=2mm,right=2mm,top=1.5mm,bottom=1.5mm,before skip=2.5mm,after skip=2.5mm,title={#1},fonttitle=\bfseries,coltitle=white}
\newtcolorbox{hinweis}{enhanced,breakable,colback=yellow!12,colframe=orange!70!black,boxrule=0.6pt,arc=2mm,left=2mm,right=2mm,top=1mm,bottom=1mm,before skip=2mm,after skip=2mm}
\newtcolorbox{ausblick}[1][Ausblick]{enhanced,breakable,colback=teal!5,colframe=teal!60!black,boxrule=0.8pt,arc=2mm,left=2mm,right=2mm,top=1.5mm,bottom=1.5mm,before skip=2.5mm,after skip=2.5mm,title={#1},fonttitle=\bfseries,coltitle=white}
\newcommand{\web}[1]{{\small\color{gray}\textit{Website: #1}}}

% ---------- Lücken, Linien, Karos ----------
\newcommand{\luecke}[1][3.5cm]{\rule[-0.5ex]{#1}{0.4pt}}
\newcommand{\linien}[1]{\par\vspace{1.5mm}\foreach \i in {1,...,#1}{\noindent\rule[0pt]{\linewidth}{0.3pt}\par\vspace{5.5mm}}\vspace{-3mm}}
\newcommand{\kariert}[1]{\par\vspace{1.5mm}\noindent\begin{tikzpicture}\draw[step=5mm,gray!55,thin] (0,0) grid (\linewidth-1pt,#1*5mm);\end{tikzpicture}\par}
\newcommand{\karobox}[2]{\begin{tikzpicture}\draw[step=5mm,gray!55,thin] (0,0) grid (#1,#2);\end{tikzpicture}}
\newcommand{\code}[1]{\texttt{#1}}

% Lücke mit Lösung: \lsg[Mindestbreite]{Lösungstext}
\newlength{\lsgbreite}
\newcommand{\lsg}[2][3.5cm]{\ifloesung\settowidth{\lsgbreite}{\,#2\,}\ifdim\lsgbreite<#1\setlength{\lsgbreite}{#1}\fi\underline{\makebox[\lsgbreite][l]{\lsgfarbe\,#2}}\else\luecke[#1]\fi}
% Schreiblinien mit Lösung: \lsglinien{Anzahl}{Lösungstext} (gleiche Höhe wie \linien)
\newcommand{\lsglinien}[2]{\ifloesung\par\vspace{1.5mm}\noindent\begin{minipage}[t][#1\dimexpr 6.9mm\relax][t]{\linewidth}\lsgfarbe #2\end{minipage}\par\vspace{-1.5mm}\else\linien{#1}\fi}
% Beliebiger Inhalt: \lsgoder{nur in der Lösung}{nur im Blatt}
\newcommand{\lsgoder}[2]{\ifloesung#1\else#2\fi}

% ---------- Graphen zeichnen ----------
\tikzset{
  knoten/.style={draw,rounded corners=3pt,fill=orange!12,minimum height=6mm,inner sep=3pt,font=\small},
  kknoten/.style={draw,circle,fill=orange!12,minimum size=7mm,inner sep=1pt,font=\small},
  kante/.style={thick},
  gkante/.style={thick,-{Stealth[length=2.6mm]}},
  dkante/.style={thick,{Stealth[length=2.6mm]}-{Stealth[length=2.6mm]}},
  gewicht/.style={font=\scriptsize,fill=white,inner sep=1pt,text=teal!60!black},
  lsgkante/.style={thick,red!75!black},
  lsggkante/.style={thick,red!75!black,-{Stealth[length=2.6mm]}},
  lsgknoten/.style={knoten,draw=red!75!black,text=red!75!black,fill=red!4},
}
% Adjazenzmatrix zum Ausfüllen: \adjmx[Lösungseinträge]{Anzahl}{Knotenliste}{Zellbreite}
% Lösungseinträge als Zeile/Spalte/Text, z. B. \adjmx[1/2/X,2/1/X]{3}{A,B,C}{7mm}
\newcommand{\adjmx}[4][]{\begin{tikzpicture}[x=#4,y=#4]
  \foreach \n [count=\i] in {#3} {\node[font=\small\bfseries] at (\i,0) {\n}; \node[font=\small\bfseries,anchor=east] at (0.45,-\i) {\n};}
  \foreach \k in {0,...,#2} {\draw[gray] (\k+0.5,-0.5) -- (\k+0.5,-#2-0.5); \draw[gray] (0.5,-\k-0.5) -- (#2+0.5,-\k-0.5);}
  \ifloesung\ifstrempty{#1}{}{\foreach \z/\s/\t in {#1} {\node[text=red!75!black,font=\small] at (\s,-\z) {\t};}}\fi
\end{tikzpicture}}

% ---------- Java-Listings ----------
\lstset{language=Java,basicstyle=\ttfamily\small,keywordstyle=\bfseries,commentstyle=\color{gray!80!black},stringstyle=\color{black},showstringspaces=false,columns=flexible,keepspaces=true,tabsize=3,frame=single,framerule=0.4pt,rulecolor=\color{black},xleftmargin=2mm,framexleftmargin=1mm,aboveskip=2mm,belowskip=2mm,breaklines=true,escapechar=§,
 literate={ä}{{\"a}}1 {ö}{{\"o}}1 {ü}{{\"u}}1 {Ä}{{\"A}}1 {Ö}{{\"O}}1 {Ü}{{\"U}}1 {ß}{{\ss}}1 {–}{{--}}1 {„}{{\glqq}}1 {“}{{\grqq}}1}
\lstnewenvironment{java}[1][]{\lstset{#1}}{}

%  Übersetzen mit XeLaTeX, LuaLaTeX, pdfLaTeX oder Tectonic.
%
%  Blatt:   xelatex ab-3-1-grundbegriffe.tex
%  Lösung:  xelatex -jobname=ab-3-1-grundbegriffe-lsg "\def\mitloesung{}\documentclass[11pt,a4paper]{article}
% =====================================================================
%  Gemeinsame Vorlage der Arbeitsblätter "Graphen" – Informatik 11 (NTG)
%  (gleiches Layout wie informatik/10/java/arbeitsblaetter/ab-vorlage.tex)
%  Einbinden in jedem Blatt direkt nach \documentclass: \input{ab-vorlage}
%  Übersetzen mit XeLaTeX, LuaLaTeX, pdfLaTeX oder Tectonic.
%
%  Blatt:   xelatex ab-3-1-grundbegriffe.tex
%  Lösung:  xelatex -jobname=ab-3-1-grundbegriffe-lsg "\def\mitloesung{}\input{ab-3-1-grundbegriffe}"
%           (füllt alle Lücken rot aus; siehe \lsg, \lsglinien, \lsgoder)
% =====================================================================
\usepackage[a4paper,left=18mm,right=18mm,top=26mm,bottom=17mm,headheight=15mm,headsep=4mm,footskip=8mm]{geometry}
\usepackage{iftex}
\ifPDFTeX
  \usepackage[T1]{fontenc}
  \usepackage[utf8]{inputenc}
  \usepackage{helvet}
  \usepackage[scaled=0.86]{beramono}
  \renewcommand{\familydefault}{\sfdefault}
\else
  \usepackage{fontspec}
  \setmainfont{texgyreheros}[Extension=.otf,UprightFont=*-regular,BoldFont=*-bold,ItalicFont=*-italic,BoldItalicFont=*-bolditalic]
  \setmonofont{DejaVuSansMono}[Extension=.ttf,UprightFont=*,BoldFont=*-Bold,ItalicFont=*-Oblique,BoldItalicFont=*-BoldOblique,Scale=0.86]
\fi
\usepackage[ngerman,shorthands=off]{babel}
\usepackage{xcolor,graphicx,tabularx,array,enumitem,fancyhdr,tikz,listings,multicol,calc,etoolbox}
\usepackage[most]{tcolorbox}
\usetikzlibrary{arrows.meta,positioning,calc,decorations.pathreplacing,shapes.geometric}
\usepackage[hidelinks]{hyperref}
\setlength{\parindent}{0pt}
\setlength{\parskip}{3pt}
\setlist{nosep,leftmargin=*}

% ---------- Lösungsschalter ----------
\newif\ifloesung
\ifdefined\mitloesung\loesungtrue\fi
\newcommand{\lsgfarbe}{\color{red!75!black}}

% ---------- Kopf- und Fußzeile (einheitliche Form) ----------
\newcommand{\lehrkraft}{Hau}
\newcommand{\themenbereich}{3. Graphen}
\newcommand{\abnummer}{}
\pagestyle{fancy}\fancyhf{}
\renewcommand{\headrulewidth}{0.4pt}
\fancyhead[L]{\small Klasse 11\\Datum:}
\fancyhead[C]{\small Themenbereich:\\\themenbereich}
\fancyhead[R]{\small \lehrkraft\ifloesung\\{\lsgfarbe\bfseries Lösung}\fi}
\fancyfoot[C]{\scriptsize edu-mrh.de\,·\,Informatik 11\,·\,Blatt \abnummer\ifloesung\ (Lösung)\fi\,·\,Seite \thepage}
\newcommand{\abtitel}[2]{\renewcommand{\abnummer}{#1}\begin{center}\Large\bfseries #1\quad\underline{#2}\end{center}\vspace{-1mm}}

% ---------- Boxen ----------
\newenvironment{aufgabe}[2][]{\begin{tcolorbox}[enhanced,breakable,colback=white,colframe=black,boxrule=0.7pt,arc=3mm,left=2mm,right=2mm,top=1.5mm,bottom=1.5mm,before skip=2.5mm,after skip=2.5mm]\textbf{Aufgabe #2)}\ifstrempty{#1}{}{ \textit{#1}}\par\vspace{0.6mm}}{\end{tcolorbox}}
\newtcolorbox{merke}[1][Merke]{enhanced,breakable,colback=green!5,colframe=green!40!black,boxrule=0.8pt,arc=2mm,left=2mm,right=2mm,top=1.5mm,bottom=1.5mm,before skip=2.5mm,after skip=2.5mm,title={#1},fonttitle=\bfseries,coltitle=white}
\newtcolorbox{hinweis}{enhanced,breakable,colback=yellow!12,colframe=orange!70!black,boxrule=0.6pt,arc=2mm,left=2mm,right=2mm,top=1mm,bottom=1mm,before skip=2mm,after skip=2mm}
\newtcolorbox{ausblick}[1][Ausblick]{enhanced,breakable,colback=teal!5,colframe=teal!60!black,boxrule=0.8pt,arc=2mm,left=2mm,right=2mm,top=1.5mm,bottom=1.5mm,before skip=2.5mm,after skip=2.5mm,title={#1},fonttitle=\bfseries,coltitle=white}
\newcommand{\web}[1]{{\small\color{gray}\textit{Website: #1}}}

% ---------- Lücken, Linien, Karos ----------
\newcommand{\luecke}[1][3.5cm]{\rule[-0.5ex]{#1}{0.4pt}}
\newcommand{\linien}[1]{\par\vspace{1.5mm}\foreach \i in {1,...,#1}{\noindent\rule[0pt]{\linewidth}{0.3pt}\par\vspace{5.5mm}}\vspace{-3mm}}
\newcommand{\kariert}[1]{\par\vspace{1.5mm}\noindent\begin{tikzpicture}\draw[step=5mm,gray!55,thin] (0,0) grid (\linewidth-1pt,#1*5mm);\end{tikzpicture}\par}
\newcommand{\karobox}[2]{\begin{tikzpicture}\draw[step=5mm,gray!55,thin] (0,0) grid (#1,#2);\end{tikzpicture}}
\newcommand{\code}[1]{\texttt{#1}}

% Lücke mit Lösung: \lsg[Mindestbreite]{Lösungstext}
\newlength{\lsgbreite}
\newcommand{\lsg}[2][3.5cm]{\ifloesung\settowidth{\lsgbreite}{\,#2\,}\ifdim\lsgbreite<#1\setlength{\lsgbreite}{#1}\fi\underline{\makebox[\lsgbreite][l]{\lsgfarbe\,#2}}\else\luecke[#1]\fi}
% Schreiblinien mit Lösung: \lsglinien{Anzahl}{Lösungstext} (gleiche Höhe wie \linien)
\newcommand{\lsglinien}[2]{\ifloesung\par\vspace{1.5mm}\noindent\begin{minipage}[t][#1\dimexpr 6.9mm\relax][t]{\linewidth}\lsgfarbe #2\end{minipage}\par\vspace{-1.5mm}\else\linien{#1}\fi}
% Beliebiger Inhalt: \lsgoder{nur in der Lösung}{nur im Blatt}
\newcommand{\lsgoder}[2]{\ifloesung#1\else#2\fi}

% ---------- Graphen zeichnen ----------
\tikzset{
  knoten/.style={draw,rounded corners=3pt,fill=orange!12,minimum height=6mm,inner sep=3pt,font=\small},
  kknoten/.style={draw,circle,fill=orange!12,minimum size=7mm,inner sep=1pt,font=\small},
  kante/.style={thick},
  gkante/.style={thick,-{Stealth[length=2.6mm]}},
  dkante/.style={thick,{Stealth[length=2.6mm]}-{Stealth[length=2.6mm]}},
  gewicht/.style={font=\scriptsize,fill=white,inner sep=1pt,text=teal!60!black},
  lsgkante/.style={thick,red!75!black},
  lsggkante/.style={thick,red!75!black,-{Stealth[length=2.6mm]}},
  lsgknoten/.style={knoten,draw=red!75!black,text=red!75!black,fill=red!4},
}
% Adjazenzmatrix zum Ausfüllen: \adjmx[Lösungseinträge]{Anzahl}{Knotenliste}{Zellbreite}
% Lösungseinträge als Zeile/Spalte/Text, z. B. \adjmx[1/2/X,2/1/X]{3}{A,B,C}{7mm}
\newcommand{\adjmx}[4][]{\begin{tikzpicture}[x=#4,y=#4]
  \foreach \n [count=\i] in {#3} {\node[font=\small\bfseries] at (\i,0) {\n}; \node[font=\small\bfseries,anchor=east] at (0.45,-\i) {\n};}
  \foreach \k in {0,...,#2} {\draw[gray] (\k+0.5,-0.5) -- (\k+0.5,-#2-0.5); \draw[gray] (0.5,-\k-0.5) -- (#2+0.5,-\k-0.5);}
  \ifloesung\ifstrempty{#1}{}{\foreach \z/\s/\t in {#1} {\node[text=red!75!black,font=\small] at (\s,-\z) {\t};}}\fi
\end{tikzpicture}}

% ---------- Java-Listings ----------
\lstset{language=Java,basicstyle=\ttfamily\small,keywordstyle=\bfseries,commentstyle=\color{gray!80!black},stringstyle=\color{black},showstringspaces=false,columns=flexible,keepspaces=true,tabsize=3,frame=single,framerule=0.4pt,rulecolor=\color{black},xleftmargin=2mm,framexleftmargin=1mm,aboveskip=2mm,belowskip=2mm,breaklines=true,escapechar=§,
 literate={ä}{{\"a}}1 {ö}{{\"o}}1 {ü}{{\"u}}1 {Ä}{{\"A}}1 {Ö}{{\"O}}1 {Ü}{{\"U}}1 {ß}{{\ss}}1 {–}{{--}}1 {„}{{\glqq}}1 {“}{{\grqq}}1}
\lstnewenvironment{java}[1][]{\lstset{#1}}{}

% Dörfer-Graph (Merkkasten Begriffe): Koordinaten wie im Begriffe-Explorer der Website
\newcommand{\doerferknoten}{%
  \node[knoten,font=\tiny,minimum height=4mm,inner sep=1.5pt] (A) at (4.4,3.3) {Altdorf};
  \node[knoten,font=\tiny,minimum height=4mm,inner sep=1.5pt] (W) at (1.2,2.9) {Weiler};
  \node[knoten,font=\tiny,minimum height=4mm,inner sep=1.5pt] (F) at (5.1,2.0) {Fischbach};
  \node[knoten,font=\tiny,minimum height=4mm,inner sep=1.5pt] (Z) at (2.9,2.0) {Ziegelstein};
  \node[knoten,font=\tiny,minimum height=4mm,inner sep=1.5pt] (N) at (0.9,0.95) {Neustadt};
  \node[knoten,font=\tiny,minimum height=4mm,inner sep=1.5pt] (B) at (4.4,0.8) {Burg};
  \node[knoten,font=\tiny,minimum height=4mm,inner sep=1.5pt] (R) at (2.6,0.2) {Rain};}
\begin{document}
\abtitel{3.1}{Graphen -- Grundbegriffe}

Hänsel und Gretel haben die Hexe überlistet. Doch die Brotkrumen, mit denen sie ihren Weg markiert hatten, haben die Vögel gefressen. Findest du trotzdem nach Hause?

\begin{aufgabe}[Hänsel und Gretel]{1}
Starte das Spiel auf der Website (Station 3.1).
\textbf{a)} Zeichne beim Spielen eine Karte des Waldes: jeder Ort ein Kästchen, jeder Weg eine Linie. Führt ein Weg nur in eine Richtung, zeichne einen Pfeil. Das Hexenhaus ist schon eingezeichnet.
\end{aufgabe}
\noindent\begin{tikzpicture}
\draw[step=5mm,gray!55,thin] (0,0) grid (17.4,7);
\node[knoten,fill=green!15] (hex) at (4,2.3) {Hexenhaus (Start)};
\ifloesung
  \node[lsgknoten] (warn) at (5.5,6.5) {Warnung};
  \node[lsgknoten] (wolf) at (10.5,6.5) {Werwolf};
  \node[lsgknoten] (stadt) at (1.6,5.2) {Stadt};
  \node[lsgknoten] (abz) at (5.5,5.2) {Abzweigung};
  \node[lsgknoten] (spinne) at (13,5.2) {Spinne};
  \node[lsgknoten] (baeck) at (1.6,3.6) {Bäckerei};
  \node[lsgknoten] (pilze) at (4,3.9) {Pilze};
  \node[lsgknoten] (w3) at (13,3.4) {Wald3};
  \node[lsgknoten] (w1) at (8,1.0) {Wald1};
  \node[lsgknoten] (stein) at (12.2,2.0) {Steinkreis};
  \node[lsgknoten] (w2) at (15,0.8) {Wald2};
  \node[lsgknoten] (heim) at (16.2,2.6) {Zuhause};
  \draw[lsgkante] (stadt) -- (baeck);
  \draw[lsgkante] (stadt) -- (pilze);
  \draw[lsgkante] (pilze) -- (abz);
  \draw[lsgkante] (pilze) -- (hex);
  \draw[lsggkante] (pilze) to[bend left=50] node[right,font=\scriptsize,text=red!75!black] {Pilze essen} (hex);
  \draw[lsgkante] (abz) -- (warn);
  \draw[lsggkante] (warn) -- (wolf);
  \draw[lsgkante] (abz) -- (spinne);
  \draw[lsgkante] (spinne) -- (w3);
  \draw[lsgkante] (hex) -- (w1);
  \draw[lsgkante] (w1) -- (w3);
  \draw[lsgkante] (w1) -- (w2);
  \draw[lsgkante] (w1) -- (stein);
  \draw[lsgkante] (w2) -- (w3);
  \draw[lsgkante] (w2) -- (stein);
  \draw[lsgkante] (w3) -- (stein);
  \draw[lsggkante] (w2) -- (heim);
  \node[text=red!75!black,font=\scriptsize,align=left,anchor=north west] at (0.2,1.3) {13 Knoten, 17 Kanten\\(14 ungerichtet, 3 gerichtet)};
\fi
\end{tikzpicture}

\textbf{b)} Am Ende zeigt dir das Spiel deinen Weg. Notiere ihn und die Anzahl der Schritte:
\lsglinien{1}{individuell; kürzester Weg: Hexenhaus -- Wald1 -- Wald2 -- Zuhause (3 Schritte)}

\begin{minipage}[t]{0.6\linewidth}
\begin{merke}[Graph und Kurzschreibweise $G = (V, E)$]\raggedright
Ein Graph besteht aus \lsg[2.3cm]{Knoten} (engl. \textit{vertices}) und \lsg[2.3cm]{Kanten} (engl. \textit{edges}). Eine Kante verbindet zwei Knoten. Im Spiel sind die Orte die \lsg[2cm]{Knoten}, die Wege die \lsg[2cm]{Kanten}.\\[1mm]
Kurzschreibweise $G = (V, E)$: $V$ ist die Menge der Knoten, $E$ die Menge der Kanten.
\begin{itemize}
\item $(A, B)$: \lsg[2.5cm]{gerichtete} Kante, nur von $A$ nach $B$
\item $\{A, B\}$: \lsg[2.5cm]{ungerichtete} Kante, steht für $(A, B)$ und \lsg[1.4cm]{$(B, A)$}
\end{itemize}
\end{merke}
\end{minipage}\hfill
\begin{minipage}[t]{0.37\linewidth}
\vspace{1mm}
\centering\small \textbf{Beispiel: Wer folgt wem?}\\[1mm]
\begin{tikzpicture}[x=0.95cm,y=0.9cm]
\node[knoten] (leon) at (0.7,2.2) {Leon};
\node[knoten] (letifa) at (3.3,2.4) {Letifa};
\node[knoten] (luise) at (4.3,1.1) {Luise};
\node[knoten] (hamza) at (2.4,0.3) {Hamza};
\node[knoten] (sabine) at (0.5,0.6) {Sabine};
\draw[dkante] (leon) -- (letifa);
\draw[gkante] (letifa) -- (luise);
\draw[gkante] (leon) -- (luise);
\draw[dkante] (leon) -- (hamza);
\draw[gkante] (hamza) -- (luise);
\end{tikzpicture}
\end{minipage}\\[1mm]
{\small $V = \{\text{Hamza}, \text{Leon}, \text{Letifa}, \text{Luise}, \text{Sabine}\}$\\
$E = \{\{\text{Leon}, \text{Letifa}\}, \{\text{Leon}, \text{Hamza}\}, (\text{Letifa}, \text{Luise}), (\text{Leon}, \text{Luise}), (\text{Hamza}, \text{Luise})\}$}

\begin{aufgabe}[Hänsel und Gretel]{1}
\textbf{c)} Schreibe den Ausschnitt deiner Karte mit den Knoten Hexenhaus, Pilze, Stadt, Bäckerei und Wald1 in Kurzschreibweise auf (5 Kanten).\\[2mm]
$V = \{$\lsg[15.4cm]{Hexenhaus, Pilze, Stadt, Bäckerei, Wald1}$\}$\\[3mm]
$E = \{$\lsg[15.4cm]{\{Hexenhaus, Pilze\}, (Pilze, Hexenhaus), \{Pilze, Stadt\}, \{Stadt, Bäckerei\}, \{Hexenhaus, Wald1\}}$\}$
\end{aufgabe}

\newpage
\begin{minipage}[t]{0.71\linewidth}
\begin{merke}[Pfade, Zyklen, Richtung, Gewicht]\raggedright\small
\textbf{Pfad (Weg):} Folge von Knoten, in der aufeinanderfolgende Knoten durch eine \lsg[1.8cm]{Kante} verbunden sind; keine Kante wird doppelt benutzt.\\[0.6mm]
\textbf{Einfacher Pfad:} Pfad, in dem sich kein \lsg[1.8cm]{Knoten} wiederholt.\\
Beispiel in A: \lsg[5.6cm]{Altdorf -- Weiler -- Ziegelstein -- Neustadt}\\[0.6mm]
\textbf{Länge eines Pfades:} ungewichtet die \lsg[3.2cm]{Anzahl der Kanten}, gewichtet die \lsg[3.2cm]{Summe der Gewichte}.\\[0.6mm]
\textbf{Zyklus:} einfacher Pfad, dessen \lsg[5cm]{Start- und Endknoten gleich sind}.\\
Beispiel in A: \lsg[5.6cm]{Weiler -- Ziegelstein -- Fischbach -- Weiler}\\
Graph mit Zyklus: \emph{zyklisch}, sonst \lsg[3.4cm]{azyklisch (zyklenfrei)}.\\[0.6mm]
\textbf{Gerichtet:} Kanten nur \lsg[3cm]{in Pfeilrichtung} (Pfeil).\ \textbf{Ungerichtet:} \lsg[2.8cm]{in beide Richtungen} (Strich).\\[0.6mm]
\textbf{Stark zusammenhängend:} Von \lsg[2.4cm]{jedem Knoten} gibt es einen Pfad zu \lsg[3.2cm]{jedem anderen Knoten}.\\
\textbf{Schwach zusammenhängend:} \lsg[5.4cm]{nur ohne Beachtung der Richtungen}\\ \lsg[\linewidth]{hängt alles zusammen. B ist schwach, denn von Rain führt keine Kante weg.}\\[0.6mm]
\textbf{Gewichtet:} Jede Kante trägt eine \lsg[2.6cm]{Zahl (Gewicht)}, z.\,B. \lsg[3.6cm]{Entfernung, Fahrzeit}.
\end{merke}
\end{minipage}\hfill
\begin{minipage}[t]{0.27\linewidth}
\vspace{1mm}
\textbf{\small A} {\scriptsize ungerichtet}\\
\begin{tikzpicture}[x=0.8cm,y=0.85cm]
\doerferknoten
\draw[kante] (A)--(W); \draw[kante] (F)--(A); \draw[kante] (W)--(F); \draw[kante] (W)--(Z);
\draw[kante] (Z)--(F); \draw[kante] (Z)--(N); \draw[kante] (B)--(Z); \draw[kante] (F)--(B);
\draw[kante] (N)--(R); \draw[kante] (B)--(R);
\end{tikzpicture}\\[2mm]
\textbf{\small B} {\scriptsize gerichtet}\\
\begin{tikzpicture}[x=0.8cm,y=0.85cm]
\doerferknoten
\draw[gkante] (A)--(W); \draw[gkante] (F)--(A); \draw[gkante] (W)--(F); \draw[gkante] (W)--(Z);
\draw[gkante] (Z)--(F); \draw[gkante] (Z)--(N); \draw[gkante] (B)--(Z); \draw[gkante] (F)--(B);
\draw[gkante] (N)--(R); \draw[gkante] (B)--(R);
\end{tikzpicture}
\end{minipage}

\begin{aufgabe}{2}
{\raggedright Zeichne $G = (V, E)$ mit $V = \{A, B, C, D, E, F\}$ und\\ $E = \{(A, B), (A, C), (A, E), (B, C), (C, C), (C, D), (E, C), (F, D)\}$. Beantworte die Fragen.\par}\vspace{1mm}
\begin{minipage}[t]{0.42\linewidth}
\vspace{0pt}
\begin{tikzpicture}
\draw[step=5mm,gray!55,thin] (0,0) grid (7,3);
\ifloesung
  \foreach \n/\x/\y/\t in {a/0.6/1.5/A, b/2/2.55/B, e/2/0.45/E, c/3.4/1.5/C, d/4.9/1.5/D, f/6.4/1.5/F} {
    \node[kknoten,minimum size=6mm,draw=red!75!black,text=red!75!black] (\n) at (\x,\y) {\t};}
  \draw[lsggkante] (a)--(b); \draw[lsggkante] (a)--(c); \draw[lsggkante] (a)--(e);
  \draw[lsggkante] (b)--(c); \draw[lsggkante] (c)--(d); \draw[lsggkante] (e)--(c);
  \draw[lsggkante] (f)--(d);
  \draw[lsggkante] (c) to[out=60,in=120,looseness=6] (c);
\fi
\end{tikzpicture}
\end{minipage}\hfill
\begin{minipage}[t]{0.55\linewidth}
\vspace{0pt}\small\raggedright
Gerichtet? Woran erkennbar? \lsg[3.9cm]{ja, Paare ( ) statt \{ \}}\\[1.5mm]
Ein Zyklus: \lsg[5.7cm]{(C, C), eine Schlinge (Länge 1)}\\[1.5mm]
Stark/schwach zusammenhängend? \lsg[3.6cm]{schwach: von D}\\[1.5mm]
\lsg[0.98\linewidth]{kommt man nirgends hin.}\\[1.5mm]
Länge des längsten einfachen Pfades: \lsg[3.2cm]{3 (A--B--C--D)}
\end{minipage}
\end{aufgabe}

\begin{aufgabe}[Zurück in den Wald]{3}
Deine Karte aus Aufgabe 1: \textbf{a)} gerichtet?\ \textbf{b)} ein Zyklus?\ \textbf{c)} stark oder schwach zusammenhängend?\ \textbf{d)} „Der Weg zurück … scheint verschwunden“: Ist die Kante Hexenhaus -- Wald1 gerichtet? Prüfe im Spiel.
\lsglinien{2}{\small a) gemischt; gerichtet nur (Pilze, Hexenhaus), (Warnung, Werwolf), (Wald2, Zuhause)\quad b) Wald1 -- Wald2 -- Wald3 -- Wald1\\ c) schwach: Von Zuhause/Werwolf führt keine Kante weg.\quad d) nein: Von Wald2 aus führt ein Weg zurück.}
\end{aufgabe}

\begin{aufgabe}[Ordne die Graphen ein.]{4}
\vspace{-1mm}
\small\renewcommand{\arraystretch}{1.15}
\begin{tabularx}{\linewidth}{|>{\centering\arraybackslash}m{5.6cm}|>{\centering\arraybackslash}m{1.7cm}|>{\centering\arraybackslash}m{1.1cm}|>{\centering\arraybackslash}m{1.1cm}|>{\centering\arraybackslash}m{1.1cm}|>{\centering\arraybackslash}X|}\hline
\textbf{Graph} & \scriptsize\textbf{zusammen-\newline hängend?}\newline stark/schwach/nein & \scriptsize\textbf{ge-\newline richtet?} & \scriptsize\textbf{ge-\newline wichtet?} & \scriptsize\textbf{zyklisch?} & \scriptsize\textbf{Länge eines längsten einfachen Pfades}\\\hline
\rule{0pt}{8mm}\textbf{I}\ \begin{tikzpicture}[baseline=(k5.base),x=0.85cm,y=0.55cm]
\foreach \n/\x/\y in {1/0/1,2/1/1,3/2/1,4/3/1,5/1/0,6/3/0} {\node[kknoten,minimum size=5mm,inner sep=0pt,font=\tiny] (k\n) at (\x,\y) {K\n};}
\draw[kante] (k1)--(k2)--(k3)--(k4); \draw[kante] (k2)--(k5);
\end{tikzpicture}
& \lsgoder{\lsgfarbe nein}{} & \lsgoder{\lsgfarbe nein}{} & \lsgoder{\lsgfarbe nein}{} & \lsgoder{\lsgfarbe nein}{} & \lsgoder{\lsgfarbe 3: K1--K2--K3--K4}{}\\\hline
\rule{0pt}{5mm}\textbf{II}\ $V = \{A, B, C\}$\newline $E =\{(A, B), (B, A), (C, B), (C, C)\}$
& \lsgoder{\lsgfarbe schwach}{} & \lsgoder{\lsgfarbe ja}{} & \lsgoder{\lsgfarbe nein}{} & \lsgoder{\lsgfarbe ja}{} & \lsgoder{\lsgfarbe 2: C$\to$B$\to$A}{}\\\hline
\rule{0pt}{8mm}\textbf{III}\ \begin{tikzpicture}[baseline=(k3.base),x=0.9cm,y=0.6cm]
\foreach \n/\x/\y in {1/0/1,2/2/1,3/1/0,4/3/0} {\node[kknoten,minimum size=5mm,inner sep=0pt,font=\tiny] (k\n) at (\x,\y) {K\n};}
\draw[kante] (k1) -- node[gewicht] {4} (k2); \draw[kante] (k2) -- node[gewicht] {1} (k3);
\draw[kante] (k3) -- node[gewicht] {5} (k1); \draw[kante] (k3) -- node[gewicht] {2} (k4);
\end{tikzpicture}
& \lsgoder{\lsgfarbe stark}{} & \lsgoder{\lsgfarbe nein}{} & \lsgoder{\lsgfarbe ja}{} & \lsgoder{\lsgfarbe ja}{} & \lsgoder{\lsgfarbe 11: K2--K1--K3--K4}{}\\\hline
\end{tabularx}
\end{aufgabe}
\end{document}
"
%           (füllt alle Lücken rot aus; siehe \lsg, \lsglinien, \lsgoder)
% =====================================================================
\usepackage[a4paper,left=18mm,right=18mm,top=26mm,bottom=17mm,headheight=15mm,headsep=4mm,footskip=8mm]{geometry}
\usepackage{iftex}
\ifPDFTeX
  \usepackage[T1]{fontenc}
  \usepackage[utf8]{inputenc}
  \usepackage{helvet}
  \usepackage[scaled=0.86]{beramono}
  \renewcommand{\familydefault}{\sfdefault}
\else
  \usepackage{fontspec}
  \setmainfont{texgyreheros}[Extension=.otf,UprightFont=*-regular,BoldFont=*-bold,ItalicFont=*-italic,BoldItalicFont=*-bolditalic]
  \setmonofont{DejaVuSansMono}[Extension=.ttf,UprightFont=*,BoldFont=*-Bold,ItalicFont=*-Oblique,BoldItalicFont=*-BoldOblique,Scale=0.86]
\fi
\usepackage[ngerman,shorthands=off]{babel}
\usepackage{xcolor,graphicx,tabularx,array,enumitem,fancyhdr,tikz,listings,multicol,calc,etoolbox}
\usepackage[most]{tcolorbox}
\usetikzlibrary{arrows.meta,positioning,calc,decorations.pathreplacing,shapes.geometric}
\usepackage[hidelinks]{hyperref}
\setlength{\parindent}{0pt}
\setlength{\parskip}{3pt}
\setlist{nosep,leftmargin=*}

% ---------- Lösungsschalter ----------
\newif\ifloesung
\ifdefined\mitloesung\loesungtrue\fi
\newcommand{\lsgfarbe}{\color{red!75!black}}

% ---------- Kopf- und Fußzeile (einheitliche Form) ----------
\newcommand{\lehrkraft}{Hau}
\newcommand{\themenbereich}{3. Graphen}
\newcommand{\abnummer}{}
\pagestyle{fancy}\fancyhf{}
\renewcommand{\headrulewidth}{0.4pt}
\fancyhead[L]{\small Klasse 11\\Datum:}
\fancyhead[C]{\small Themenbereich:\\\themenbereich}
\fancyhead[R]{\small \lehrkraft\ifloesung\\{\lsgfarbe\bfseries Lösung}\fi}
\fancyfoot[C]{\scriptsize edu-mrh.de\,·\,Informatik 11\,·\,Blatt \abnummer\ifloesung\ (Lösung)\fi\,·\,Seite \thepage}
\newcommand{\abtitel}[2]{\renewcommand{\abnummer}{#1}\begin{center}\Large\bfseries #1\quad\underline{#2}\end{center}\vspace{-1mm}}

% ---------- Boxen ----------
\newenvironment{aufgabe}[2][]{\begin{tcolorbox}[enhanced,breakable,colback=white,colframe=black,boxrule=0.7pt,arc=3mm,left=2mm,right=2mm,top=1.5mm,bottom=1.5mm,before skip=2.5mm,after skip=2.5mm]\textbf{Aufgabe #2)}\ifstrempty{#1}{}{ \textit{#1}}\par\vspace{0.6mm}}{\end{tcolorbox}}
\newtcolorbox{merke}[1][Merke]{enhanced,breakable,colback=green!5,colframe=green!40!black,boxrule=0.8pt,arc=2mm,left=2mm,right=2mm,top=1.5mm,bottom=1.5mm,before skip=2.5mm,after skip=2.5mm,title={#1},fonttitle=\bfseries,coltitle=white}
\newtcolorbox{hinweis}{enhanced,breakable,colback=yellow!12,colframe=orange!70!black,boxrule=0.6pt,arc=2mm,left=2mm,right=2mm,top=1mm,bottom=1mm,before skip=2mm,after skip=2mm}
\newtcolorbox{ausblick}[1][Ausblick]{enhanced,breakable,colback=teal!5,colframe=teal!60!black,boxrule=0.8pt,arc=2mm,left=2mm,right=2mm,top=1.5mm,bottom=1.5mm,before skip=2.5mm,after skip=2.5mm,title={#1},fonttitle=\bfseries,coltitle=white}
\newcommand{\web}[1]{{\small\color{gray}\textit{Website: #1}}}

% ---------- Lücken, Linien, Karos ----------
\newcommand{\luecke}[1][3.5cm]{\rule[-0.5ex]{#1}{0.4pt}}
\newcommand{\linien}[1]{\par\vspace{1.5mm}\foreach \i in {1,...,#1}{\noindent\rule[0pt]{\linewidth}{0.3pt}\par\vspace{5.5mm}}\vspace{-3mm}}
\newcommand{\kariert}[1]{\par\vspace{1.5mm}\noindent\begin{tikzpicture}\draw[step=5mm,gray!55,thin] (0,0) grid (\linewidth-1pt,#1*5mm);\end{tikzpicture}\par}
\newcommand{\karobox}[2]{\begin{tikzpicture}\draw[step=5mm,gray!55,thin] (0,0) grid (#1,#2);\end{tikzpicture}}
\newcommand{\code}[1]{\texttt{#1}}

% Lücke mit Lösung: \lsg[Mindestbreite]{Lösungstext}
\newlength{\lsgbreite}
\newcommand{\lsg}[2][3.5cm]{\ifloesung\settowidth{\lsgbreite}{\,#2\,}\ifdim\lsgbreite<#1\setlength{\lsgbreite}{#1}\fi\underline{\makebox[\lsgbreite][l]{\lsgfarbe\,#2}}\else\luecke[#1]\fi}
% Schreiblinien mit Lösung: \lsglinien{Anzahl}{Lösungstext} (gleiche Höhe wie \linien)
\newcommand{\lsglinien}[2]{\ifloesung\par\vspace{1.5mm}\noindent\begin{minipage}[t][#1\dimexpr 6.9mm\relax][t]{\linewidth}\lsgfarbe #2\end{minipage}\par\vspace{-1.5mm}\else\linien{#1}\fi}
% Beliebiger Inhalt: \lsgoder{nur in der Lösung}{nur im Blatt}
\newcommand{\lsgoder}[2]{\ifloesung#1\else#2\fi}

% ---------- Graphen zeichnen ----------
\tikzset{
  knoten/.style={draw,rounded corners=3pt,fill=orange!12,minimum height=6mm,inner sep=3pt,font=\small},
  kknoten/.style={draw,circle,fill=orange!12,minimum size=7mm,inner sep=1pt,font=\small},
  kante/.style={thick},
  gkante/.style={thick,-{Stealth[length=2.6mm]}},
  dkante/.style={thick,{Stealth[length=2.6mm]}-{Stealth[length=2.6mm]}},
  gewicht/.style={font=\scriptsize,fill=white,inner sep=1pt,text=teal!60!black},
  lsgkante/.style={thick,red!75!black},
  lsggkante/.style={thick,red!75!black,-{Stealth[length=2.6mm]}},
  lsgknoten/.style={knoten,draw=red!75!black,text=red!75!black,fill=red!4},
}
% Adjazenzmatrix zum Ausfüllen: \adjmx[Lösungseinträge]{Anzahl}{Knotenliste}{Zellbreite}
% Lösungseinträge als Zeile/Spalte/Text, z. B. \adjmx[1/2/X,2/1/X]{3}{A,B,C}{7mm}
\newcommand{\adjmx}[4][]{\begin{tikzpicture}[x=#4,y=#4]
  \foreach \n [count=\i] in {#3} {\node[font=\small\bfseries] at (\i,0) {\n}; \node[font=\small\bfseries,anchor=east] at (0.45,-\i) {\n};}
  \foreach \k in {0,...,#2} {\draw[gray] (\k+0.5,-0.5) -- (\k+0.5,-#2-0.5); \draw[gray] (0.5,-\k-0.5) -- (#2+0.5,-\k-0.5);}
  \ifloesung\ifstrempty{#1}{}{\foreach \z/\s/\t in {#1} {\node[text=red!75!black,font=\small] at (\s,-\z) {\t};}}\fi
\end{tikzpicture}}

% ---------- Java-Listings ----------
\lstset{language=Java,basicstyle=\ttfamily\small,keywordstyle=\bfseries,commentstyle=\color{gray!80!black},stringstyle=\color{black},showstringspaces=false,columns=flexible,keepspaces=true,tabsize=3,frame=single,framerule=0.4pt,rulecolor=\color{black},xleftmargin=2mm,framexleftmargin=1mm,aboveskip=2mm,belowskip=2mm,breaklines=true,escapechar=§,
 literate={ä}{{\"a}}1 {ö}{{\"o}}1 {ü}{{\"u}}1 {Ä}{{\"A}}1 {Ö}{{\"O}}1 {Ü}{{\"U}}1 {ß}{{\ss}}1 {–}{{--}}1 {„}{{\glqq}}1 {“}{{\grqq}}1}
\lstnewenvironment{java}[1][]{\lstset{#1}}{}

\begin{document}
\abtitel{3.3}{Breitensuche}

\begin{aufgabe}[Silvesterfeuerwerk]{1}
\begin{minipage}[t]{0.45\linewidth}
\vspace{0pt}\raggedright
Die Stadt hat Feuerwerksbatterien mit Zündschnüren verbunden. Batterie A wird um 0\,s gezündet, jede Zündschnur brennt genau 1\,s. Welche Batterien zünden wann?\\[2mm]
\renewcommand{\arraystretch}{1.35}
\begin{tabular}{|c|p{4.3cm}|}\hline
\textbf{Zeit} & \textbf{Batterien}\\\hline
0\,s & A\\\hline
1\,s & \lsgoder{\lsgfarbe B, C, D, E}{}\\\hline
2\,s & \lsgoder{\lsgfarbe F, G, H}{}\\\hline
3\,s & \lsgoder{\lsgfarbe I, J (zuletzt)}{}\\\hline
\end{tabular}
\end{minipage}\hfill
\begin{minipage}[t]{0.5\linewidth}
\vspace{0pt}\centering
\begin{tikzpicture}[x=1.25cm,y=0.95cm]
\foreach \n/\x/\y in {I/0/3, C/1.7/3.3, B/3.6/3.3, F/0.9/2.3, A/2.7/2.3, D/4.3/1.8, E/2.0/1.2, H/1.0/0.3, G/3.0/0.3, J/2.0/-0.5} {
  \node[kknoten,minimum size=6mm] (\n) at (\x,\y) {\n};}
\node[font=\scriptsize,anchor=west] at (A.east) {\,Start};
\foreach \a/\b in {A/B, A/C, A/D, A/E, B/D, C/F, E/F, E/H, E/G, F/I, H/J, G/J} {\draw[kante] (\a)--(\b);}
\end{tikzpicture}
\end{minipage}
\end{aufgabe}

\begin{merke}[Breitensuche]\raggedright
Die Breitensuche durchläuft (\emph{traversiert}) einen Graphen \lsg[2.8cm]{schichtweise}: zuerst alle Nachbarn des Startknotens, dann deren noch nicht besuchte Nachbarn usw. Mit jeder Schicht wird die Pfadlänge um \lsg[1.8cm]{eins} größer. Deshalb findet sie in ungewichteten Graphen den Pfad mit den \lsg[3.2cm]{wenigsten Kanten}.\\[1mm]
Die noch zu besuchenden Knoten stehen in einer \lsg[3.2cm]{Warteschlange}: Wer zuerst hineinkommt, kommt zuerst wieder heraus (engl. \emph{first in, first out}). Besuchte Knoten werden \lsg[2.8cm]{markiert}, damit kein Knoten doppelt besucht wird -- wie die Brotkrumen im Spiel.
\end{merke}

\begin{aufgabe}{2}
\begin{minipage}[t]{0.4\linewidth}
\vspace{0pt}\raggedright
\textbf{a)} Führe die Breitensuche ab Knoten A durch. Nimm Nachbarn immer alphabetisch. Notiere nach jedem Schritt den aktuellen Knoten und die Warteschlange.\\[2mm]
\begin{tikzpicture}[x=1.05cm,y=0.95cm]
\foreach \n/\x/\y in {A/2.6/3, B/0.3/2.3, C/3.6/2.2, G/4.7/2.9, H/5.7/2.0, F/1.9/1.0, D/4.6/1.1, E/3.2/0.5, I/4.3/-0.1} {
  \node[kknoten,minimum size=6mm] (\n) at (\x,\y) {\n};}
\foreach \a/\b in {A/B, A/C, A/F, C/G, C/H, B/D, D/H, D/E, E/F, E/I} {\draw[kante] (\a)--(\b);}
\end{tikzpicture}
\end{minipage}\hfill
\begin{minipage}[t]{0.56\linewidth}
\vspace{0pt}\small
\renewcommand{\arraystretch}{1.18}
\begin{tabularx}{\linewidth}{|c|c|X|}\hline
\textbf{Schritt} & \textbf{aktueller Knoten} & \textbf{Warteschlange danach}\\\hline
0 & -- & A\\\hline
1 & \lsgoder{\lsgfarbe A}{} & \lsgoder{\lsgfarbe B C F}{}\\\hline
2 & \lsgoder{\lsgfarbe B}{} & \lsgoder{\lsgfarbe C F D}{}\\\hline
3 & \lsgoder{\lsgfarbe C}{} & \lsgoder{\lsgfarbe F D G H}{}\\\hline
4 & \lsgoder{\lsgfarbe F}{} & \lsgoder{\lsgfarbe D G H E}{}\\\hline
5 & \lsgoder{\lsgfarbe D}{} & \lsgoder{\lsgfarbe G H E}{}\\\hline
6 & \lsgoder{\lsgfarbe G}{} & \lsgoder{\lsgfarbe H E}{}\\\hline
7 & \lsgoder{\lsgfarbe H}{} & \lsgoder{\lsgfarbe E}{}\\\hline
8 & \lsgoder{\lsgfarbe E}{} & \lsgoder{\lsgfarbe I}{}\\\hline
9 & \lsgoder{\lsgfarbe I}{} & \lsgoder{\lsgfarbe (leer)}{}\\\hline
\end{tabularx}
\end{minipage}\\[2mm]
\begin{minipage}[t]{0.44\linewidth}
\vspace{0pt}\raggedright
\textbf{b)} Zeichne den \emph{Besuchsbaum}: einen Pfeil von jedem Knoten zu den Knoten, die er in die Warteschlange gestellt hat.\\[1mm]
\begin{tikzpicture}
\draw[step=5mm,gray!55,thin] (0,0) grid (7,3);
\ifloesung
  \foreach \n/\x/\y in {A/3.5/2.6, B/1.2/1.6, C/3.5/1.6, F/5.8/1.6, D/0.6/0.45, G/2.8/0.45, H/4.2/0.45, E/5.8/0.45, I/6.7/0.45} {
    \node[kknoten,minimum size=5mm,inner sep=0pt,font=\scriptsize,draw=red!75!black,text=red!75!black] (\n) at (\x,\y) {\n};}
  \foreach \a/\b in {A/B, A/C, A/F, B/D, C/G, C/H, F/E, E/I} {\draw[lsggkante] (\a)--(\b);}
\fi
\end{tikzpicture}
\end{minipage}\hfill
\begin{minipage}[t]{0.53\linewidth}
\vspace{0pt}\raggedright\small
\textbf{c)} Reihenfolge ab Knoten E:\\[1mm]
\lsg[\linewidth]{E, D, F, I, B, H, A, C, G}\\[2mm]
\textbf{d)} Anton behauptet: „Die Breitensuche besucht immer alle Knoten.“ Nimm Stellung.\\[1mm]
\lsg[\linewidth]{Falsch: Nur Knoten, die vom Start aus erreichbar sind.}\\[1mm]
\lsg[\linewidth]{Bei nicht zusammenhängenden oder nur schwach}\\[1mm]
\lsg[\linewidth]{zusammenhängenden Graphen fehlen Knoten.}
\end{minipage}
\end{aufgabe}

\newpage
\begin{aufgabe}[Projekt Breitensuche]{3}
\textbf{a)} Ordne jeder Zeile des Pseudocodes die Nummer des passenden Codefragments zu (Website bzw. Datei \code{Codefragmente.java}).\quad
\textbf{b)} Setze die Methode \code{breitensuche} im Projekt um (TODO 1 bis 3) und teste sie.
\end{aufgabe}
{\renewcommand{\arraystretch}{1.25}\small
\begin{tabularx}{\linewidth}{|X|c|}\hline
\textbf{Pseudocode} \code{breitensuche(start, ziel)} & \textbf{Fragment}\\\hline
aktuellerKnoten $\leftarrow$ start & \lsgoder{\lsgfarbe 2}{}\\\hline
lege aktuellerKnoten in die Warteschlange & \lsgoder{\lsgfarbe 1}{}\\\hline
gib „Reihenfolge: “ aus & \lsgoder{\lsgfarbe 7}{}\\\hline
\textbf{solange} aktuellerKnoten $\neq$ ziel \textbf{und} Warteschlange nicht leer & \lsgoder{\lsgfarbe 6}{}\\\hline
\quad aktuellerKnoten $\leftarrow$ vordersten Knoten aus der Warteschlange nehmen & \lsgoder{\lsgfarbe 4}{}\\\hline
\quad aktuellerKnoten als besucht markieren & \lsgoder{\lsgfarbe 8}{}\\\hline
\quad gib aktuellerKnoten aus & \lsgoder{\lsgfarbe 3}{}\\\hline
\quad \textbf{für alle} Knoten $i$ & \lsgoder{\lsgfarbe 5}{}\\\hline
\quad\quad \textbf{falls} Kante von aktuellerKnoten zu $i$ \textbf{und} $i$ nicht besucht \textbf{und} nicht in der Warteschlange:\newline\quad\quad\quad stelle $i$ hinten an die Warteschlange & \lsgoder{\lsgfarbe 9}{}\\\hline
\textbf{falls} aktuellerKnoten $=$ ziel: „gefunden“, \textbf{sonst}: „nicht gefunden“ & \lsgoder{\lsgfarbe 10}{}\\\hline
\end{tabularx}}

\begin{aufgabe}[Projekt Breitensuche: Testen]{4}
\begin{minipage}[t]{0.5\linewidth}
\vspace{0pt}\raggedright
Im Hauptprogramm entsteht dieser Graph. Knoten 5 wird mit \code{knotenHinzufuegen()} ergänzt, hat aber keine Kante.\\[1mm]
\textbf{a)} Sage die Ausgabe voraus, dann teste.
\end{minipage}\hfill
\begin{minipage}[t]{0.46\linewidth}
\vspace{0pt}\centering
\begin{tikzpicture}[x=1.1cm,y=0.75cm]
\foreach \n/\x/\y in {0/0/2, 1/1.5/1.2, 2/3/2, 3/0/0.4, 4/3/0.2, 5/4.3/1.1} {\node[kknoten,minimum size=6mm] (k\n) at (\x,\y) {\n};}
\foreach \a/\b in {0/1, 2/1, 1/3, 0/3, 2/4} {\draw[kante] (k\a)--(k\b);}
\end{tikzpicture}
\end{minipage}\\[1mm]
{\renewcommand{\arraystretch}{1.35}\small
\begin{tabularx}{\linewidth}{|l|X|c|}\hline
\textbf{Aufruf} & \textbf{Reihenfolge} & \textbf{gefunden?}\\\hline
\code{breitensuche(3, 2)} & \lsgoder{\lsgfarbe 3 0 1 2}{} & \lsgoder{\lsgfarbe ja}{}\\\hline
\code{breitensuche(4, 0)} & \lsgoder{\lsgfarbe 4 2 1 0}{} & \lsgoder{\lsgfarbe ja}{}\\\hline
\code{breitensuche(0, 5)} & \lsgoder{\lsgfarbe 0 1 3 2 4}{} & \lsgoder{\lsgfarbe nein}{}\\\hline
\end{tabularx}}\\[2mm]
\textbf{b)} \textit{Stolperstein:} Verschiebe Fragment 4 (\code{poll}) an das Ende der Schleife. Was gibt \code{breitensuche(3, 2)} jetzt aus? Erkläre, was schiefgeht.
\lsglinien{3}{„Reihenfolge: 3 3 0“, nicht gefunden. Der Startknoten liegt noch in der Warteschlange und wird ein zweites Mal geholt. Nach dem Holen von Knoten 1 ist die Warteschlange leer, die Schleife endet, bevor Knoten 1 untersucht wird: Knoten 2 wird nie erreicht.}
\end{aufgabe}
\end{document}
